\documentclass[11pt]{amsproc}

\usepackage[a4paper,margin=2.6cm]{geometry}
\usepackage{amsmath,amssymb,amsthm}
\usepackage[T1]{fontenc}
\usepackage[utf8]{inputenc}

\newcommand{\C}{\mathbb{C}}
\newcommand{\tr}{\operatorname{tr}}
\newcommand{\Id}{\operatorname{Id}}

\begin{document}

\section*{Problema}

Para cada entero $n\ge 0$ sea $P_n$ el espacio vectorial de los polinomios
homogéneos de grado $n$ en las variables $x,y,z$, con coeficientes complejos.
El polinomio nulo pertenece a todos los $P_n$, y ponemos $P_n=\{0\}$ si $n<0$.
Consideremos el operador
\[
  Df=\frac{\partial^3 f}{\partial x^3}+\frac{\partial^3 f}{\partial y^3}
      +\frac{\partial^3 f}{\partial z^3},
\]
que lleva $P_n$ en $P_{n-3}$, el subespacio
\[
  V_n=\{f\in P_n : Df=0\}
\]
y la transformación lineal $T\colon P_n\to P_n$ dada por $(Tf)(x,y,z)=f(y,z,x)$.

\begin{enumerate}
  \item\itshape Probar que $D(P_n)=P_{n-3}$ para todo $n\ge 3$. Es decir, que
        todo polinomio homogéneo de grado $n-3$ es de la forma
        $g_{xxx}+g_{yyy}+g_{zzz}$ con $g\in P_n$.
  \item\itshape Probar que $T(V_n)\subseteq V_n$ y calcular, para cada $n\ge 0$,
        el polinomio característico de la restricción de $T$ a $V_n$.
\end{enumerate}

Dicho de otro modo, en (2) se pide, para cada $\lambda\in\C$ con $\lambda^3=1$,
la dimensión del espacio de soluciones $f\in P_n$ de la ecuación
$f_{xxx}+f_{yyy}+f_{zzz}=0$ que cumplen $f(y,z,x)=\lambda f(x,y,z)$.

\section*{Solución}

\subsection*{Parte (1)}

Fijemos $n\ge 3$. Todo $h\in P_{n-3}$ se escribe de manera única como
\[
  h=\sum_{k=0}^{n-3}x^{n-3-k}h_k(y,z),
\]
con $h_k$ homogéneo de grado $k$ en $y,z$. Llamemos $\kappa(h)$ al mayor $k$
con $h_k\ne 0$ y pongamos $\kappa(0)=-1$. Probemos, por inducción en
$\kappa(h)$, que $h\in D(P_n)$.

Para $\kappa(h)=-1$ tenemos $h=0=D(0)$. Supongamos ahora que $K=\kappa(h)\ge 0$
y que el resultado vale para los polinomios de $P_{n-3}$ con $\kappa<K$.
Integremos tres veces respecto de $x$ el término de mayor grado en $y,z$. Como
$K\le n-3$, el número $c=(n-K)(n-K-1)(n-K-2)$ no es nulo. Tomemos
\[
  g=\frac{1}{c}\,x^{n-K}h_K(y,z)\in P_n .
\]
Al derivar tres veces respecto de $x$ recuperamos el término que elegimos.
Las derivadas respecto de $y$ y de $z$ bajan en tres el grado en esas variables.
Así,
\[
  Dg=x^{n-3-K}h_K+\frac{1}{c}\,x^{n-K}
        \bigl(\partial_y^3h_K+\partial_z^3h_K\bigr).
\]
El segundo sumando es nulo o tiene grado $K-3$ en $y,z$. Luego
$\kappa(h-Dg)<K$ y, por la hipótesis inductiva, $h-Dg=Dg'$ para cierto
$g'\in P_n$. De aquí $h=D(g+g')$, como queríamos.

Como $\dim P_n=\binom{n+2}{2}$ para $n\ge 0$, el teorema de la dimensión da
\[
  \dim V_n=\dim P_n-\dim P_{n-3}
          =\binom{n+2}{2}-\binom{n-1}{2}=3n\qquad (n\ge 3).
\]
Para $n=0,1,2$ se tiene $V_n=P_n$, de dimensiones $1,3,6$, respectivamente.
En particular, $\dim V_n=3n$ para todo $n\ge 1$.

\subsection*{Parte (2), paso 1}

En $Tf$ la variable $x$ ocupa el tercer lugar de $f$, la $y$ el primero y
la $z$ el segundo. Por la regla de la cadena,
\[
  \partial_x(Tf)=T(\partial_z f),\qquad
  \partial_y(Tf)=T(\partial_x f),\qquad
  \partial_z(Tf)=T(\partial_y f).
\]
Aplicando estas igualdades tres veces y sumando, obtenemos $D(Tf)=T(Df)$.
En particular, si $Df=0$ entonces $D(Tf)=0$, de modo que $T(V_n)\subseteq V_n$.
Notemos también que $T^3=\Id$.

\subsection*{Parte (2), paso 2}

Calculemos la traza de $T$ en $V_n$. Probemos primero que
\begin{equation}\label{eq:traza}
  \tr(T|_{V_n})=\tr(T|_{P_n})-\tr(T|_{P_{n-3}}).
\end{equation}
Para $0\le n\le 2$ esto se sigue de que $V_n=P_n$ y $P_{n-3}=\{0\}$.
Para $n\ge 3$, tomemos una base de $V_n$ y completémosla con $g_1,\dots,g_r$
a una base de $P_n$. Los vectores $Dg_1,\dots,Dg_r$ generan $P_{n-3}$ por
la parte (1), y forman una base porque $r=\dim P_n-\dim V_n=\dim P_{n-3}$.

Escribamos $Tg_k=v_k+\sum_l\beta_{lk}g_l$, con $v_k\in V_n$. Aplicando $D$
y usando el paso 1, obtenemos
\[
  T(Dg_k)=D(Tg_k)=\sum_l\beta_{lk}\,Dg_l.
\]
Así, en la base elegida de $P_n$ la matriz de $T$ es triangular por bloques.
Sus bloques diagonales son la matriz de $T|_{V_n}$ y la matriz
$(\beta_{lk})$ de $T|_{P_{n-3}}$ en la base $(Dg_l)$. Tomando trazas se
obtiene \eqref{eq:traza}.

Ahora bien, $T$ lleva el monomio $x^ay^bz^c$ en $x^cy^az^b$, de modo que su
traza en $P_n$ es la cantidad de monomios fijos. Un monomio es fijo si y sólo
si $a=b=c$. Para $n\ge 3$, multiplicar por $xyz$ da una biyección entre los
monomios fijos de grado $n-3$ y los de grado $n$. Para $n=1,2$ tampoco hay
monomios fijos y $P_{n-3}=\{0\}$. Por \eqref{eq:traza} resulta entonces
\[
  \tr(T|_{V_n})=0\qquad(n\ge 1).
\]

\subsection*{Parte (2), paso 3}

Sea $\omega=e^{2\pi i/3}$. Como $T^3=\Id$ y el polinomio
$t^3-1=(t-1)(t-\omega)(t-\omega^2)$ tiene raíces simples, $T|_{V_n}$ es
diagonalizable. Fijemos $n\ge 1$ y llamemos $m_0,m_1,m_2$ a las
multiplicidades de $1,\omega,\omega^2$. Usando que $1+\omega+\omega^2=0$,
la igualdad del paso anterior se escribe
\[
  0=m_0+\omega m_1+\omega^2m_2=(m_0-m_2)+(m_1-m_2)\omega.
\]
Como los $m_i$ son enteros y $\omega$ no es real, la parte imaginaria da
$m_1=m_2$ y luego $m_0=m_2$. Su suma es $\dim V_n=3n$, así que cada uno
vale $n$. Para $n=0$, $V_0$ es el espacio de las constantes y $T=\Id$.
En conclusión,
\[
  \chi_{T|_{V_n}}(t)=
  \begin{cases}
    t-1,       & n=0,\\
    (t^3-1)^n, & n\ge 1.
  \end{cases}
\]

\section*{Observaciones}

\begin{enumerate}
  \item La primera parte sirve de pista. La sobreyectividad dice que el
    núcleo tiene la menor dimensión posible y sugiere comparar $P_n$ con
    $P_{n-3}$. Para la versión difícil basta con pedir el segundo inciso.
    Si se quiere dar una pista para el primero, se puede preguntar qué pasa
    con el grado en $y,z$ al integrar tres veces respecto de $x$ el término
    de mayor grado en $y,z$ y aplicar después el operador $D$.

  \item Para $n\ge 3$, $V_n$ no tiene una base de monomios. Por ejemplo,
    $(x-y)^n$ pertenece a $V_n$ porque sus dos terceras derivadas en $x$ y
    en $y$ se cancelan, mientras que $x^n$, que aparece en su desarrollo,
    no pertenece a $V_n$. En cambio, en $P_n$ las trazas de las potencias de
    $T$ se leen directamente en la base de monomios. De hecho, para $r=0,1,2$,
    \[
      m_r=\frac{1}{3}\sum_{j=0}^2\omega^{-rj}\tr(T^j|_{V_n}).
    \]
    Un autovalor $\omega^s$ aporta al miembro derecho
    $\tfrac13\sum_{j=0}^2\omega^{(s-r)j}$, que vale $1$ si $s=r$ y $0$ en
    los demás casos. El grupo simétrico $S_3$ actúa en $V_n$ permutando las
    variables. La fórmula da el producto escalar del carácter de esta
    representación, restringido al grupo cíclico generado por la permutación
    $(x,y,z)\mapsto(y,z,x)$, con el carácter que le asigna $\omega^r$ a ese
    ciclo. Es la misma cuenta que hace la función
    \textsf{EigenvaluesChar} de \textsf{GAP}.

  \item Hay otra manera de terminar el segundo inciso. Sea $W_n$ el subespacio
    generado por los monomios de grado $n$ a los que les falta alguna variable.
    Éstos son precisamente los monomios que resuelven $f_{xyz}=0$, y tenemos
    \[
      P_n=xyz\,P_{n-3}\oplus W_n .
    \]
    Ambos sumandos son estables por $T$, y multiplicar por $xyz$ conmuta con
    $T$. Por lo tanto, el polinomio característico de $T$ en $P_n$ es el
    producto de los correspondientes a $P_{n-3}$ y a $W_n$. Por la matriz
    triangular por bloques del paso 2, ese mismo polinomio es el producto de
    los correspondientes a $P_{n-3}$ y a $V_n$. Cancelando el factor común,
    resulta $\chi_{T|_{V_n}}=\chi_{T|_{W_n}}$.

    Para $n\ge 1$, los monomios de $W_n$ forman ciclos de longitud tres bajo
    $T$. Hay exactamente $n$ ciclos, representados por $x^ay^{n-a}$ con
    $1\le a\le n$. En efecto, en cada ciclo de monomios con dos variables
    hay uno solo al que le falta $z$, y el ciclo de las potencias puras está
    representado por $x^n$. Cada ciclo aporta un factor $t^3-1$.

  \item La misma respuesta vale para $D=q(\partial_x,\partial_y,\partial_z)$
    con $q$ cualquier polinomio homogéneo cúbico no nulo que cumpla
    $q(x,y,z)=q(y,z,x)$. La regla de la cadena da de nuevo $DT=TD$. Para
    ver la sobreyectividad, tomemos en cada $P_n$ el producto interno para
    el cual los monomios son ortogonales y
    $\langle x^ay^bz^c,x^ay^bz^c\rangle=a!\,b!\,c!$. Respecto de este
    producto, $\partial_x$ es la adjunta de la multiplicación por $x$, y lo
    mismo vale para $y,z$. Así, $q(\partial)$ es la adjunta de la
    multiplicación por $\bar q$, donde se conjugan los coeficientes. Como
    $\bar q\ne 0$, esa multiplicación es inyectiva de $P_{n-3}$ en $P_n$.
    Su adjunta es entonces suryectiva y el resto de la solución vale sin cambios.

  \item Con la ecuación de Laplace $f_{xx}+f_{yy}+f_{zz}=0$ se obtiene el
    espacio de polinomios armónicos homogéneos de grado $n$. La permutación
    cíclica de las coordenadas es una rotación de ángulo $2\pi/3$ alrededor
    del eje $(1,1,1)$. Toda forma cuadrática no nula invariante por $T$ es
    combinación de $x^2+y^2+z^2$ y $(x+y+z)^2$, de modo que también es
    invariante por todas las rotaciones de ese eje. En el caso cúbico esto
    ya falla para $x^3+y^3+z^3$. Sobre el plano $x+y+z=0$ coincide con
    $3xyz$, que cambia de signo al girar media vuelta alrededor del eje.

    También se puede calcular la respuesta cuadrática con el argumento
    anterior. Para cualquier forma cuadrática $q$ como la de arriba,
    $q(\partial)$ es suryectiva por el mismo argumento de adjuntas. Con la
    notación $\chi_m=\chi_{T|_{P_m}}$, el polinomio buscado es
    $\chi_n/\chi_{n-2}$. Para calcularlo, tomemos las coordenadas lineales
    $u=x+y+z$, $v=x+\omega^2y+\omega z$ y $w=x+\omega y+\omega^2z$. Sus
    autovalores son $1,\omega,\omega^2$, respectivamente, así que $vw$ es
    fijo por $T$. Como en el caso de $xyz$, basta considerar los monomios a
    los que les falta $v$ o $w$. Son $u^n$ y los pares
    $u^{n-j}v^j,u^{n-j}w^j$ con $1\le j\le n$. Sus autovalores son $1$ y
    los pares $\omega^j,\omega^{-j}$. Cada múltiplo de $3$ aporta dos unos,
    y cada uno de los demás índices aporta un $\omega$ y un $\omega^2$.
    Resulta así el polinomio
    \[
      (t-1)^{2\lfloor n/3\rfloor+1}
      (t^2+t+1)^{n-\lfloor n/3\rfloor}.
    \]
\end{enumerate}

\end{document}

